Every claim about the hardware rests on bit-exact equality with the integer reference. Table~\ref{tab:verif} summarizes
the evidence.

\begin{table}[t]
\centering\footnotesize
\caption{Verification evidence (all simulations with assertions on address ranges, INT32 accumulator overflow and the
divider's precondition).}
\label{tab:verif}
\setlength{\tabcolsep}{3pt}
\begin{tabular}{@{}p{0.43\columnwidth}p{0.5\columnwidth}@{}}
\toprule
Check & Result \\ \midrule
Every trained model, test set, all logits & \zooVerifRuns{} checkpoints, \zooVerifiedLogits{} logits, 0 mismatches \\
Same, specialized engine and load port & 0 mismatches \\
One-bitstream engine via AXI (10 models) & 0 mismatches; 200-token story identical \\
Randomized differential tests & \zooFuzzCases{} random architectures/sizes/weights, \zooFuzzLogits{} logits, 0 mismatches \\
4-state simulation (Icarus), X on outputs & 0 unknown logits \\
Divider, exhaustive ($W{=}12$) + random & \zooArithCases{} cases, 0 errors \\
Independent C implementation of the ISA & bit-exact on all models \\
On-board (PYNQ-Z2), all models & \zooBoardVerif{} \\ \bottomrule
\end{tabular}
\end{table}

The simulations use Verilator 5.052; the randomized tests draw an architecture, vocabulary (27--512), context (8--64),
depth (1--3) and width at random, fill the weights with values of random scale, calibrate on random tokens, and compare
every logit of random token sequences. The divider's skip-ahead is proven by exhaustive simulation of every valid input
at $W=12$ and by the property $q\,b\le a<(q+1)\,b$ on random 64-bit cases; the softmax reciprocal is checked for every
denominator in $[2^{15},2^{24})$. A third implementation, a C instruction-set simulator of the microcode, matches as well
and serves as the CPU baseline.
